\subsection{序向量空间}

向量空间 E 上的一个预序结构，记作 \(x \preccurlyeq y\) 或 \(y \succcurlyeq x\)，称为与 E 的向量空间结构相容，如果它满足下列两条公理：

\((\mathrm{EO}_{\mathrm{I}})\) 若 \(x \preccurlyeq y\)，则 \(x + z \preccurlyeq y + z\)（对所有 \(z \in \mathrm{E}\)）。

\((\mathrm{EO}_{\mathrm{II}})\) 若 \(x \succcurlyeq 0\)，则 \(\lambda x \succcurlyeq 0\)（对每个标量 \(\lambda \geqslant 0\)）。

赋有这两个结构的向量空间 E 称为预序向量空间（相应地，当 E 上的预序关系是一个序时，称为序向量空间）。

注意，公理 \((EO_{I})\) 意味着 E 的预序结构与加法群结构相容，也就是说，赋有这两个结构的 E 是一个预序群 (A, VI, 第 3 页)。

例. --- 在定义在 A 上的所有实值函数组成的空间 \(E = R^{A}\) 上，由「对所有 \(t \in A\)，\(x(t) \leqslant y(t)\)」给出的序关系与 E 的向量空间结构相容。

命题 13. --- (i) 预序向量空间 E 中由 \(\succcurlyeq 0\) 的元素组成的集 P 是尖凸锥。

\begin{enumerate}
\def\labelenumi{(\roman{enumi})}
\setcounter{enumi}{1}
\item
  反之，若 \(\mathbf{P}\) 是 \(\mathbf{E}\) 中的一个尖凸锥，则关系 \(y - x \in \mathbf{P}\) 是 \(\mathbf{E}\) 上的一个预序关系，且由它定义的预序结构是唯一与 E 的向量空间结构相容且使 P 为由 \(\succcurlyeq 0\) 的元素组成之集的预序结构。
\item
  关系 \(y - x \in \mathbf{P}\)（其中 \(\mathbf{P}\) 是尖凸锥）是 \(\mathbf{E}\) 上的序关系，当且仅当 \(\mathbf{P}\) 是真尖锥。
\item
  公理 \((\mathrm{EO}_{\mathrm{I}})\) 与 \((\mathrm{EO}_{\mathrm{II}})\) 蕴含 \(P + P \subset P\) 及 \(\lambda P \subset P\)（对所有 \(\lambda > 0\)）。由于 \(0 \in P\)，由此推出 P 是尖凸锥（II, 第 11 页, prop. 10）。
\item
  反之，若 \(\mathbf{P}\) 是尖凸锥，则关系 \(\mathbf{P} + \mathbf{P} \subset \mathbf{P}\) 蕴含关系 \(y - x \in \mathbf{P}\) 是与 E 的加法群结构相容的预序 (A, VI, 第 3 页, prop. 3)；显然，把它写成 \(x \preccurlyeq y\) 时，集 \(\mathbf{P}\) 与由 \(x \geqslant 0\) 组成的集恒同；进而，关系「\(\lambda \mathbf{P} \subset \mathbf{P}\)（对所有 \(\lambda \geqslant 0\)）」表明公理 \((\mathrm{EO}_{\mathrm{II}})\) 满足。(iii) 说 \(\mathbf{P}\) 是真尖锥，意味着 \(\mathbf{P} \cap (-\mathbf{P}) = \{0\}\)（II, 第 11 页, cor. 2），从而 \(y - x \in \mathbf{P}\) 是序关系。
\end{enumerate}

例. --- * 设 H 是一个实 Hilbert 空间；在 H 的连续自同态组成的向量空间 \(\mathcal{L}(\mathrm{H})\) 中，正埃尔米特自同态组成一个真尖凸锥；因此这个锥定义了 \(\mathcal{L}(\mathrm{H})\) 上一个与其向量空间结构相容的序结构，且其中关系 \(A \leqslant B\) 意指 \(B - A\) 是正埃尔米特自同态。*

对向量空间 E 中任意尖凸锥 P，集 \(P \cap (-P)\) 是 E 的一个向量子空间 H（II, 第 11 页, cor. 2）。P 在 E/H 中的典范像 \(P'\) 是凸锥，且 \(P'\) 在 E 中的原像是 P。于是 \(P' \cap (-P') = \{0\}\)，从而 \(P'\) 在 E/H 上定义一个与其向量空间结构相容的序结构。

预序向量空间 E 上的线性型 f 称为正的，如果 E 中 \(x \geqslant 0\) 蕴含 \(f(x) \geqslant 0\)。或者等价地说，如果 E 中由 \(\geqslant 0\) 的元素组成的凸锥 P 包含在由满足 \(f(x) \geqslant 0\) 的 x 组成的半空间中。显然，在 E 的对偶 E* 中，正线性型组成的集是一个尖凸锥。

